\documentclass[addpoints,twoside]{exam}
\usepackage{problemset}

\begin{document}

\FontEasyRead

\Heading
{Microeconomics}
{Consumer Behavior \& Demand - Version A}
{\textbf{Instructions: }Answer each question clearly. Show reasoning where appropriate.}

\begin{questions}

    \Problem{Explain the three key assumptions about consumer preferences: completeness, transitivity, and “more is better.”}

    \Problem{Explain why indifference curves are downward sloping. What would it mean if they were upward sloping?}

    \Problem{Explain why indifference curves are typically convex (bowed inward). Use the concept of marginal rate of substitution in your answer.}

    \Problem{A consumer has a fixed income and chooses between goods $X$ and $Y$.
        \begin{itemize}
            \item What does the budget line represent?
            \item What determines its slope?
        \end{itemize}}

    \Problem{Define marginal utility. What does it mean for marginal utility to be diminishing?}

    \Problem{Explain why the optimal consumption bundle occurs where
        \[
            MRS = \frac{P_X}{P_Y}.
        \]
        What does this condition mean intuitively?}

    \Problem{A consumer has:
        \[
            MU_X = 20,\quad P_X = 4,\quad MU_Y = 15,\quad P_Y = 5
        \]
        \begin{itemize}
            \item Is the consumer maximizing utility?
            \item If not, what adjustment should they make?
        \end{itemize}}

    \Problem{A consumer is willing to give up 3 units of $Y$ to gain 1 unit of $X$.
        \begin{itemize}
            \item What is the marginal rate of substitution of $X$ for $Y$?
            \item If $\frac{P_X}{P_Y} = 2$, is the consumer at an optimal bundle? Explain.
        \end{itemize}}

    \Problem{Explain the difference between the substitution effect and the income effect when the price of a good falls.}

    \Problem{Define the following types of goods and explain how demand changes with income:
        \begin{itemize}
            \item Normal good
            \item Inferior good
        \end{itemize}}

    \Problem{What is an Engel curve? What does its slope tell you about a good?}

    \Problem{Explain how an individual demand curve is derived from a price-consumption curve.}

    \Problem{Suppose Consumer A demands 5 units of a good at a price of \$10, and Consumer B demands 7 units at the same price.
        \begin{itemize}
            \item What is total market demand at \$10?
            \item How is a market demand curve constructed from individual demand curves?
        \end{itemize}}

    \Problem{Explain why the price elasticity of demand is typically negative.}

    \Problem{Explain why elasticity varies along a linear demand curve. Why is demand more elastic at high prices?}

    \Problem{A decrease in price leads to a decrease in quantity demanded.
        \begin{itemize}
            \item What type of good could this be?
            \item What must be true about the income and substitution effects?
        \end{itemize}}

    \Problem{Define consumer surplus. How is it represented on a demand graph?}

\end{questions}

\end{document}